% ============================================================
% CHAPTER 7: IMPLEMENTATION & ARCHITECTURE
% MCDA 5585 / 5586 Major Project Report
% Bhavik Kantilal Bhagat | A00494758 | Saint Mary's University
% ============================================================

\chapter{Implementation \& Architecture}
\label{chap:implementation}

\section{Platform Overview: Four-Layer Architecture}
\label{sec:platform_overview}
% ────────────────────────────────────────────────────────────

The three deliverables of my internship form a coherent four-layer
observability stack. Figure~\ref{fig:platform_overview} shows how the layers
relate to each other and to the underlying IA platform infrastructure.

\begin{figure}[ht]
\centering
\begin{lstlisting}[language={}, basicstyle=\ttfamily\footnotesize, frame=none, backgroundcolor=\color{white}, numbers=none]
 +-----------------------------------------------------------------+
 |                    CONSUMERS / OPERATORS                        |
 |   Grafana AMG Dashboard  .  Inventory REST API  .  SRE Team    |
 +--------------------------------+--------------------------------+
                                  |
 +--------------------------------v--------------------------------+
 |              VISUALIZATION & ALERTING LAYER                    |
 |  Amazon Managed Grafana (AMG)        CloudWatch Alarms         |
 |  14 rows . 110+ panels . variables   38 alarms . SNS routing   |
 +--------------------------------+--------------------------------+
                                  |
 +--------------------------------v--------------------------------+
 |                        DATA LAYER                              |
 |  Amazon CloudWatch  |  CloudWatch Logs Insights  |  X-Ray     |
 |  Metric Insights SQL|  Saved queries (30x7 plat.)|  (PoC)     |
 +--------------------------------+--------------------------------+
                                  |
 +--------------------------------v--------------------------------+
 |      INFRASTRUCTURE LAYER (5 Budget Code Environments)         |
 |  Lambda (119) . ECS (22 clusters) . EC2 (83) . MSK/Kafka       |
 |  SQS . DocumentDB . Neptune . RDS . Redis . OpenSearch          |
 +-----------------------------------------------------------------+
\end{lstlisting}
\caption{IA-SRE Observability Platform --- High-Level Architecture (four layers)}
\label{fig:platform_overview}
\end{figure}

The \textbf{infrastructure layer} is the monitored subject: five budget
platforms housing Lambda functions, ECS clusters, EC2
instances, and the full spectrum of AWS database and messaging services
used by the Meridian and Citco Works platforms.

The \textbf{data layer} ingests, stores, and queries observability signals
from the infrastructure layer. Amazon CloudWatch is the sole data source ---
a deliberate architectural decision (see Section~\ref{sec:decision_cloudwatch_only})
that eliminates the operational overhead of a self-managed Prometheus stack.
CloudWatch Metric Insights SQL provides the per-resource aggregation capability
that drives multi-service panels without requiring one panel per resource.

The \textbf{visualization and alerting layer} is split between Amazon Managed
Grafana (AMG), which renders the 14-row dashboard for human operators, and
CloudWatch Alarms with SNS routing, which fires machine-readable notifications
when monitored metrics breach their SLO-aligned thresholds.

The \textbf{consumers layer} includes the operators who use the Grafana
dashboard during incident response, the IA Resources Inventory REST API
(IISS-507) that exposes the AWS service fleet programmatically, and the SRE
team's email distribution list that receives SNS alarm notifications.

% ────────────────────────────────────────────────────────────
\section{CloudWatch Dashboard Platform (IISS-455/469/482--486)}
\label{sec:cw_dashboard_platform}
% ────────────────────────────────────────────────────────────

\subsection{Initial Lambda Dashboard (IISS-455)}

My earliest monitoring work, completed in April~2026, produced a native
Amazon CloudWatch dashboard covering the \texttt{rpacfs1} and
\texttt{ia-bpo} Lambda function fleets --- a combined inventory of
\textbf{51~functions} at the time of first publication. The dashboard
provided per-function panels for invocation counts, error rates, duration
averages, and throttle events, and I accompanied it with a Confluence
inventory page mapping each function to its owning team and SLO target.

The architectural pattern for this early work was straightforward:
CloudWatch native widgets queried the \texttt{AWS/Lambda} namespace using
\texttt{SEARCH()} expressions to enumerate all functions matching a
\texttt{FunctionName} prefix. This approach worked adequately for small,
homogeneous fleets but surfaced two structural limitations that would
drive the later Metric Insights SQL migration:

\begin{itemize}
    \item \textbf{SEARCH() visibility gap}: The \texttt{SEARCH()} function
          only returns metrics for resources with recent data points.
          Lambda functions that have not been invoked within the lookback
          window become invisible to the dashboard, creating a false
          impression of a smaller fleet.

    \item \textbf{No cross-environment aggregation}: A single
          \texttt{SEARCH()} expression could not span multiple tag
          dimensions simultaneously, making it difficult to aggregate
          across the six platforms (\texttt{rpacfs1},
          \texttt{ia-bpo}, \texttt{automationhub}, \texttt{meridian},
          \texttt{citcoworks}) without duplicating panels.
\end{itemize}

\subsection{Extended Dashboard and Architecture Analysis (IISS-469)}

IISS-469 extended the initial dashboard with additional tag dimensions
covering \texttt{cfsbpm}, \texttt{citcoworks}, and \texttt{meridian}
budget codes, and introduced CloudWatch Logs Insights widgets for error
log surfacing. Rather than relying solely on \texttt{Errors} metrics from
the \texttt{AWS/Lambda} namespace, the Logs Insights widgets queried the
\texttt{/aws/lambda/} log group prefix directly, enabling pattern-matched
error extraction regardless of whether Lambda had published a
corresponding metric.

A significant architectural investigation within IISS-469 evaluated the
feasibility of a fully centralized log aggregation approach using
CloudWatch metric subscription filters. My analysis quantified the
infrastructure cost of this design: routing log events from all IA Lambda
functions into a centralized processing layer would require \textbf{106
metric subscription filters}---one per function per metric
dimension---across the six platforms. While technically
feasible from a CloudWatch service quota perspective, this approach was
ultimately blocked by an IAM permission boundary constraint:
the \texttt{logs:PutSubscriptionFilter} action was not granted within the
IA team's scoped permission set, preventing programmatic creation of
subscription filters at the scale required.

This investigation also evaluated five distinct approaches to out-of-memory
(OOM) detection for Lambda functions: (1) parsing \texttt{REPORT}
log lines from CloudWatch Logs Insights; (2) comparing
\texttt{MaxMemoryUsed} against \texttt{MemorySize} via the Lambda
\texttt{enhanced monitoring} extension; (3) CloudWatch Lambda Insights
(paid tier); (4) custom metric emission on memory threshold breach; and
(5) structured log parsing via a metric filter pattern. I documented each approach
with its detection latency, cost, and permission requirements.
This analysis informed the Container Insights cost evaluation I conducted
under IISS-483 and the alarm design in IISS-487.

\subsection{IA Infrastructure Monitoring Sub-Stories (IISS-482 to IISS-486)}

\subsubsection{EC2 Inventory and Tag Mismatch Discovery (IISS-483)}

My gap analysis of the EC2 inventory scanned \textbf{83~instances} across
the six platforms. This scan revealed a critical tagging
inconsistency: the \texttt{rpacfs1} environment's EC2 instances were
tagged with \texttt{BudgetCode=``CFS/RPA''} rather than the expected
\texttt{``rpacfs1''} value, making them invisible to tag-based
\texttt{BudgetCode} lookups. This discovery led me to introduce the
\texttt{AppName} fallback pattern: when a primary \texttt{BudgetCode}
lookup returns no results for \texttt{rpacfs1}, the discovery logic falls
back to filtering by \texttt{AppName=``rpacfs1''}. This fallback became
the standard for all resource discovery in the inventory API (see
Section~\ref{sec:inventory_api}).

\subsubsection{Container Insights Cost Analysis (IISS-484)}

With ECS task-level telemetry absent from five of the six Meridian ECS
clusters, IISS-484 evaluated four Container Insights deployment options:

\begin{enumerate}
    \item \textbf{No Container Insights} --- zero cost, zero task-level visibility.
    \item \textbf{Enhanced Container Insights} --- full task and pod metrics, but
          carries per-metric ingestion cost that scales with task count.
    \item \textbf{Standard Container Insights} --- cluster- and service-level
          CPU/memory metrics via the \texttt{ECS/ContainerInsights}
          namespace, charged at CloudWatch standard metric pricing.
    \item \textbf{CloudWatch Agent on ECS (self-managed)} --- maximum flexibility,
          highest operational overhead.
\end{enumerate}

My analysis recommended \textbf{Standard Container Insights} at an estimated
cost of \textbf{\$10.44/month} for full coverage across all 22~ECS clusters.
This option provided the SLI data required for FR2 (ECS CPU/memory dashboard)
at negligible incremental cost relative to the existing CloudWatch spend.
I enabled Standard Container Insights on five of the six Meridian clusters
on May~11--12, 2026.

\subsubsection{SQS Registry Architecture (IISS-482 SQS Feature)}

The SQS monitoring feature I added on June~12, 2026 introduced a
\textbf{registry-based architecture with a three-tier fallback model} for
resilient SQS queue data retrieval. The three tiers operated as follows:

\begin{enumerate}
    \item \textbf{Tier 1 --- Live API}: Query CloudWatch
          \texttt{AWS/SQS} metrics directly via the CloudWatch API. Fast
          and fresh, but requires network connectivity and API quota.

    \item \textbf{Tier 2 --- S3 Cache}: If the API call fails (rate
          limit, transient error), retrieve the last-known queue snapshot
          from an S3 bucket. Stale by up to one cache TTL window, but
          always available.

    \item \textbf{Tier 3 --- Local File}: If both API and S3 fail, fall
          back to a locally bundled queue inventory file. Guaranteed to
          return data; snapshot reflects deployment-time inventory.
\end{enumerate}

This three-tier model ensured the SQS monitoring pipeline remained
functional even during AWS service degradation events. I accompanied the feature with \textbf{960~automated tests} covering the registry
dispatch logic, fallback transitions, and queue metadata serialization.

% ────────────────────────────────────────────────────────────
\section{Grafana AMG Observability Platform (IISS-487)}
\label{sec:grafana_amg_platform}
% ────────────────────────────────────────────────────────────

IISS-487 was my primary deliverable, consuming
\textbf{140~hours} over May~6--June~18, 2026. It delivered a centralized
Amazon Managed Grafana workspace with 14~dashboard rows, 110+~panels,
38~CloudWatch alarms, and a Lambda-backed IaC provisioner --- all built
on a four-layer architecture described below.

\subsection{Architectural Layers}

\subsubsection{Layer 1: Data Layer}

Amazon CloudWatch serves as the \textbf{sole data source} for the entire
Grafana workspace. My deliberate choice to use only the native Grafana
CloudWatch plugin --- rather than a self-managed Prometheus/CloudWatch
exporter stack --- eliminated all additional infrastructure requirements.
The data layer is characterized by:

\begin{itemize}
    \item \textbf{CloudWatch plugin}: The Grafana AMG workspace connects
          to CloudWatch via the built-in plugin, authenticated through an
          IAM role. No additional data collectors, agents, or scrapers
          are required.

    \item \textbf{IAM role model}: A \textbf{16-role agent permission
          model} provides the Grafana service account with cross-account
          read access across all IA AWS accounts. Each role grants the
          minimum permissions required: \texttt{cloudwatch:GetMetricData},
          \texttt{cloudwatch:ListMetrics}, \texttt{cloudwatch:DescribeAlarms},
          \texttt{logs:StartQuery}, \texttt{logs:GetQueryResults}, and
          \texttt{ec2:DescribeTags} (for tag-based resource enumeration).

    \item \textbf{Metric Insights SQL}: The primary query language for
          multi-resource panels. CloudWatch Metric Insights enables
          SQL-like aggregation across all resources in a namespace with
          a single query, replacing the brittle \texttt{SEARCH()} expressions
          used in the earlier CloudWatch dashboard.
\end{itemize}

\subsubsection{Layer 2: Visualization Layer}

The Grafana AMG workspace hosts \textbf{14~dashboard rows} organized by
monitored service category. Table~\ref{tab:dashboard_rows} summarizes the
dashboard structure.

\begin{table}[ht]
\centering
\caption{Grafana AMG Dashboard --- Row Structure (14 Rows)}
\label{tab:dashboard_rows}
\begin{tabular}{p{0.5cm}p{3.0cm}p{8.5cm}}
\toprule
\textbf{Row} & \textbf{Service} & \textbf{Panels and Metrics} \\
\midrule
1 & SQS & Queue depth (ApproximateNumberOfMessagesVisible), oldest message age (ApproximateAgeOfOldestMessage) per queue; DLQ depth timeseries \\
\midrule
2--3 & Lambda & Invocations, errors, duration P50/P95, throttles, concurrent executions --- per budget code \\
\midrule
4--5 & ECS & CPU and memory utilisation per cluster and service; Container Insights task-level metrics (\texttt{ECS/ContainerInsights} namespace) \\
\midrule
6 & MSK/Kafka & Consumer lag: \texttt{EstimatedMaxTimeLag} and \texttt{SumOffsetLag} per consumer group \\
\midrule
7--8 & EC2 & CPU utilisation, memory (CloudWatch agent), disk I/O, network in/out --- tag-based (\texttt{AppName}) \\
\midrule
9 & DocumentDB / Neptune & Database connections, query latency \\
\midrule
10 & OpenSearch & Index rate, search latency, JVM heap pressure \\
\midrule
11 & Grafana AMG health & Workspace CPU, active users, query latency \\
\midrule
12 & Alarm state overview & Composite alarm status board (OK/ALARM/INSUFFICIENT\_DATA per platform) \\
\midrule
13 & Resource inventory & Panel linking to the IA Resources Inventory API \\
\midrule
14 & Deep-link row & Clickable AWS Console URLs for each resource \\
\bottomrule
\end{tabular}
\end{table}

\subsubsection{Grafana Template Variables}

The dashboard supports environment-level filtering through three Grafana
template variables backed by CloudWatch dimension-value queries:

\begin{itemize}
    \item \texttt{\$budget\_code} --- filters all panels to a specific
          environment (\texttt{rpacfs1}, \texttt{ia-bpo}, \texttt{meridian},
          \texttt{citcoworks}, \texttt{automationhub}). Backed by a
          CloudWatch \texttt{ListMetrics} call enumerating distinct
          \texttt{BudgetCode} tag values.

    \item \texttt{\$ecs\_cluster} --- filters ECS panels to a specific
          cluster. Backed by a \texttt{dimension\_values()} query on
          the \texttt{ECS/ContainerInsights} namespace.

    \item \texttt{\$lambda\_function} --- filters Lambda panels to a
          specific function. Backed by a \texttt{dimension\_values()}
          query on the \texttt{AWS/Lambda} namespace.
\end{itemize}

When a user selects a value from any variable, all panels that reference
that variable re-execute their queries with the selected filter applied,
effectively scoping the entire dashboard to a single environment or
resource. This eliminates the need to maintain separate dashboards per
platform.

\subsubsection{CloudWatch Metric Insights SQL: Architecture Pattern}

CloudWatch Metric Insights SQL is the key architectural pattern that
enables per-resource drill-down across large fleets within a single query.
The standard form used throughout the dashboard is:

\begin{lstlisting}[language=SQL, caption={Metric Insights SQL --- Per-Function Lambda Duration Average}]
SELECT AVG(Duration)
FROM "AWS/Lambda"
GROUP BY FunctionName
ORDER BY AVG() DESC
LIMIT 10
\end{lstlisting}

This single query returns the average duration for each Lambda function
as a separate time series, covering all Lambda functions across all
codes without requiring one panel per function. The \texttt{WHERE}
clause can filter by tag or dimension:

\begin{lstlisting}[language=SQL, caption={Metric Insights SQL --- Platform-Filtered ECS CPU}]
SELECT AVG(CpuUtilized)
FROM "ECS/ContainerInsights"
WHERE ClusterName LIKE 'meridian-%'
GROUP BY ServiceName
\end{lstlisting}

\paragraph{Percentile Limitation and Workaround.}
A fundamental constraint of Metric Insights SQL is that the
\texttt{GROUP BY} mode only supports the aggregate functions
\texttt{AVG}, \texttt{SUM}, \texttt{MIN}, and \texttt{MAX}. Percentile
statistics (\texttt{p50}, \texttt{p95}, \texttt{p99}) are \emph{not
available} in \texttt{GROUP BY} mode. This is a CloudWatch API limitation,
not a Grafana limitation.

For panels requiring true latency percentiles, the workaround is to use
the \emph{standard CloudWatch metrics} query mode with a wildcard
dimension value:

\begin{lstlisting}[language={}, basicstyle=\ttfamily\small, caption={Grafana CloudWatch Panel --- Wildcard Dimension for p99}]
{
  "metricName": "Duration",
  "namespace": "AWS/Lambda",
  "statistics": ["p99"],
  "dimensions": { "FunctionName": ["*"] }
}
\end{lstlisting}

The wildcard \texttt{FunctionName: ["*"]} returns individual time series
for each function that has emitted \texttt{Duration} metrics within the
selected time window. The tradeoff is that only recently active functions
appear --- the same ``visibility gap'' limitation noted for \texttt{SEARCH()}.
For the Lambda duration P50/P95 panels, I accepted this tradeoff:
a function that has not been invoked recently has no meaningful latency
distribution to display.

\subsubsection{Layer 3: Alerting Layer}

I provisioned \textbf{38~CloudWatch alarms} across all monitored
platforms, each connected to an SNS topic that routes notifications to
the IA-SRE team's email distribution list. Alarm categories include:

\begin{itemize}
    \item \textbf{Lambda alarms}: Error rate exceeding 1\% over a
          5-minute rolling window; throttle count exceeding threshold;
          concurrent execution count approaching account limit.

    \item \textbf{ECS alarms}: CPU utilisation sustained above 80\%
          over a 5-minute evaluation period; task desired count $\neq$
          running count (task failure detection).

    \item \textbf{MSK alarms}: \texttt{EstimatedMaxTimeLag} exceeding
          10,000~messages on any consumer group. This alarm directly
          addresses the April~17 incident: it would have fired within
          minutes of the lag beginning to build, compared to the 3.5-hour
          detection gap that occurred under the reactive model.

    \item \textbf{SQS alarms}: Queue depth exceeding configurable
          threshold per queue; \texttt{ApproximateAgeOfOldestMessage}
          exceeding 1~hour (configurable per SLA).

    \item \textbf{Composite alarms}: Platform-level health signals that
          aggregate per-service alarm states into a single
          \texttt{OK}/\texttt{ALARM} indicator per budget code
          environment. Displayed on Row~12 of the dashboard.
\end{itemize}

\subsubsection{Layer 4: Infrastructure-as-Code Provisioner}

The IaC layer is a \textbf{Lambda-backed CloudFormation provisioner}
deployed via AWS SAM. Its architecture addresses a specific scaling
constraint: with 119~Lambda functions, a static CloudFormation template
defining one alarm per function would contain $119 \times (\text{alarm
types per function})$ resources, rapidly approaching the CloudFormation
500-resource-per-stack limit.

The provisioner resolves this by \emph{dynamically discovering} the
current function fleet at deploy time and synthesizing alarm resources
programmatically:

\begin{enumerate}
    \item The provisioner Lambda function is invoked by a CloudFormation
          custom resource on every \texttt{samconfig.toml} deploy.

    \item It reads parameterized CloudFormation YAML templates from an
          S3 bucket (see Section~\ref{sec:decision_secrets_s3}).

    \item For each template (Lambda alarms, ECS alarms, SQS alarms, MSK
          alarms), it calls the relevant AWS discovery API to enumerate
          all current resources in the specified platform.

    \item It synthesizes CloudFormation \texttt{AWS::CloudWatch::Alarm}
          resource blocks for each discovered resource and passes the
          generated template to CloudFormation for stack deployment.

    \item Four CloudFormation stacks were deployed on
          \textbf{May~11, 2026}: \texttt{ia-sre-lambda-alarms},
          \texttt{ia-sre-ecs-alarms}, \texttt{ia-sre-sqs-alarms}, and
          \texttt{ia-sre-msk-alarms}.
\end{enumerate}

A \texttt{samconfig.toml} with environment-specific parameter files
enables reuse across DEV, UAT, and PROD without code changes. The budget
code, alarm thresholds, and SNS topic ARN are all externalised as
SAM parameters, satisfying NFR5 (multi-environment reuse).

\subsection{Grafana Version Upgrade: Architectural Resilience Test}

On \textbf{May~13, 2026} --- two days after the initial infrastructure
deployment --- the Grafana AMG workspace was upgraded from version
\textbf{10.4 to 12.4}. This was a major version migration that
constituted an unplanned but highly informative architectural resilience
test.

The upgrade introduced breaking changes in two areas:

\begin{itemize}
    \item \textbf{CloudWatch Metric Insights SQL editor}: The query
          syntax editor in Grafana 12.4 changed the way SQL expressions
          were associated with panel targets. All SQL panels required
          migration to the updated CloudWatch query model.

    \item \textbf{Variable interpolation}: Template variable syntax
          in panel queries changed between versions, causing several
          \texttt{\$budget\_code} and \texttt{\$ecs\_cluster} references
          to fail to resolve correctly after the upgrade.
\end{itemize}

I resolved this through a systematic audit of all 110+~panels to identify
affected queries. The migration reinforced a key design principle
documented as NFR6: \textbf{use only standard CloudWatch plugin queries
and avoid AMG-version-specific syntax}. I migrated all panel queries
to the stable Grafana 12.4 CloudWatch query model, and the workspace has
remained stable through subsequent Grafana patch releases.

% ────────────────────────────────────────────────────────────
\section{IA Resources Inventory API (IISS-507)}
\label{sec:inventory_api}
% ────────────────────────────────────────────────────────────

IISS-507 delivered a \textbf{Flask REST API} that I built for programmatic discovery
of all IA platform cloud resources. I housed the project in the
\texttt{ia-sre-resources-inventory} repository and scaffolded it on
June~11, 2026; it reached production-ready state with 562~automated
tests and 99\%~code coverage.

\subsection{Technology Stack}

\begin{itemize}
    \item \textbf{Language}: Python~3, \texttt{flask[async]} for async
          route handlers
    \item \textbf{AWS SDK}: \texttt{boto3} for all resource discovery
    \item \textbf{Validation}: Pydantic response models (type-safe
          \texttt{Resource} dataclass with fields: \texttt{name},
          \texttt{type}, \texttt{budget\_code}, \texttt{tags},
          \texttt{arn})
    \item \textbf{Deployment}: AWS Lambda via \texttt{serverless-wsgi}
          adapter
    \item \textbf{Testing}: \texttt{pytest} + \texttt{Hypothesis}
          (property-based tests)
\end{itemize}

\subsection{Discoverer Registry Pattern}

My core architectural contribution in IISS-507 is the
\textbf{Discoverer Registry pattern}. Rather than routing resource
discovery logic through a monolithic handler, the API maintains a
registry that maps \texttt{(resource\_type, budget\_code)} tuples to
concrete discoverer classes:

\begin{lstlisting}[language=Python, caption={Discoverer Registry --- Core Pattern}]
from abc import ABC, abstractmethod
from typing import List

class BaseDiscoverer(ABC):
    """Abstract base class for all resource discoverers."""

    @abstractmethod
    async def discover(self) -> List[Resource]:
        """Discover all resources of this type for the budget code."""
        ...

# Registry maps (resource_type, budget_code) -> discoverer class
REGISTRY: dict = {
    ("lambda",  "rpacfs1"):      RPACFSLambdaDiscoverer,
    ("lambda",  "ia-bpo"):       IABPOLambdaDiscoverer,
    ("ecs",     "meridian"):     MeridianECSDiscoverer,
    ("ecs",     "citcoworks"):   CitcoWorksECSDiscoverer,
    ("ec2",     "rpacfs1"):      RPACFSEc2Discoverer,
    # ... 7 more stub discoverers
}

async def discover_resources(
    resource_type: str,
    budget_code: str | None = None
) -> List[Resource]:
    """Route discovery request to the appropriate discoverer."""
    results = []
    for (rtype, bc), discoverer_cls in REGISTRY.items():
        if rtype != resource_type:
            continue
        if budget_code and bc != budget_code:
            continue
        discoverer = discoverer_cls()
        results.extend(await discoverer.discover())
    return results
\end{lstlisting}

Each discoverer implements the single abstract method
\texttt{async def discover() -> List[Resource]}. Adding support for a
new platform or resource type requires implementing one class and
registering one entry --- zero changes to routing, serialization, or
test infrastructure.

\subsection{API Endpoints}

\begin{table}[ht]
\centering
\caption{IA Resources Inventory API --- Endpoint Summary}
\label{tab:api_endpoints}
\begin{tabular}{p{5.5cm}p{7.0cm}}
\toprule
\textbf{Endpoint} & \textbf{Description} \\
\midrule
\texttt{GET /api/v1/lambda} & All Lambda functions; optional
  \texttt{?budget\_code=rpacfs1} filter \\
\midrule
\texttt{GET /api/v1/ecs} & All ECS clusters across all platforms \\
\midrule
\texttt{GET /api/v1/ec2} & All EC2 instances (tag-based discovery) \\
\midrule
\texttt{GET /api/v1/secrets} & Secrets Manager secrets inventory \\
\midrule
\texttt{GET /api/v1/s3} & S3 buckets per budget code \\
\midrule
\texttt{GET /api/v1/health} & Health check --- returns 200 OK \\
\bottomrule
\end{tabular}
\end{table}

\subsection{Tag-Based Discovery and the Two-Stage Filter}

The most technically interesting discovery challenge was the
\texttt{CITCOWORKS} and \texttt{MERIDIAN} tag ambiguity. Both
environments share the AWS tag value \texttt{``Shared Cloud Artifacts''}
for the \texttt{BudgetCode} dimension --- a legacy of the two platforms
being onboarded before the IA team established per-platform
\texttt{BudgetCode} uniqueness requirements. A naive
\texttt{BudgetCode} lookup would conflate resources from two operationally
distinct environments.

My solution is a \textbf{two-stage filter} implemented in both the
\texttt{CitcoWorksECSDiscoverer} and \texttt{MeridianECSDiscoverer}:

\begin{lstlisting}[language=Python,
  caption={Two-Stage Filter --- Disambiguating Shared BudgetCode}]
class MeridianECSDiscoverer(BaseDiscoverer):
    """Discovers ECS resources for the Meridian environment."""

    async def discover(self) -> List[Resource]:
        client = boto3.client("ecs")
        all_clusters = []

        # Stage 1: filter by BudgetCode (shared with CITCOWORKS)
        paginator = client.get_paginator("list_clusters")
        async for page in paginator.paginate():
            for arn in page["clusterArns"]:
                tags = await self._get_tags(arn)
                if tags.get("BudgetCode") != "Shared Cloud Artifacts":
                    continue
                # Stage 2: AppName discriminates Meridian vs CitcoWorks
                if tags.get("AppName", "").lower() != "meridian":
                    continue
                all_clusters.append(self._build_resource(arn, tags))

        return all_clusters
\end{lstlisting}

The \texttt{AppName} tag is now designated the \textbf{mandatory secondary
discriminator} for all resources in the \texttt{``Shared Cloud Artifacts''}
BudgetCode environments. Any new resource added to either environment
without the correct \texttt{AppName} tag will be invisible to the
inventory API. This risk is mitigated by a tag compliance CloudWatch
alarm that fires if new ECS services or Lambda functions lack the
required tags.

\subsection{Async Architecture}

The API uses \texttt{flask[async]} with async route handlers.
Each handler aggregates boto3 calls concurrently using
\texttt{asyncio.gather()}, reducing total API latency by parallelizing
the AWS API calls within a single request:

\begin{lstlisting}[language=Python,
  caption={Async Route Handler --- Concurrent Discovery}]
@app.route("/api/v1/lambda")
async def get_lambda_resources():
    budget_code = request.args.get("budget_code")
    # All discoverers run concurrently, not sequentially
    results = await asyncio.gather(*[
        discoverer_cls().discover()
        for (rtype, bc), discoverer_cls in REGISTRY.items()
        if rtype == "lambda"
        and (budget_code is None or bc == budget_code)
    ])
    all_resources = [r for sublist in results for r in sublist]
    return jsonify([r.dict() for r in all_resources])
\end{lstlisting}

This pattern reduces the API response time for a full fleet enumeration
from the sequential sum of all boto3 call latencies to the latency of
the single slowest call --- consistently keeping the API within the
10-second NFR3 target.

\subsection{Test Strategy}

My test suite for IISS-507 comprises \textbf{562~automated tests}
achieving \textbf{99\%~code coverage}:

\begin{itemize}
    \item \textbf{Unit tests}: Each discoverer class tested in isolation
          with mocked boto3 clients. Verifies correct tag filtering,
          pagination handling, and \texttt{Resource} object construction.

    \item \textbf{Integration tests}: Full Flask test client hitting
          all REST endpoints. Verifies routing, serialization, and
          error response format.

    \item \textbf{Property-based tests (Hypothesis)}: Tag-parsing logic,
          registry dispatch, and Pydantic response serialization tested
          with randomly generated inputs. Verifies that the two-stage
          filter never returns a \texttt{MERIDIAN} resource when
          \texttt{budget\_code=citcoworks} is specified, regardless of
          what combination of tag values is present.

    \item \textbf{Edge case tests}: Malformed tags, empty budget code
          environments, boto3 pagination edge cases (empty pages,
          single-page results, maximum page sizes).
\end{itemize}

% ────────────────────────────────────────────────────────────
\section{Design Decisions \& Tradeoffs}
\label{sec:design_decisions}
% ────────────────────────────────────────────────────────────

Four design decisions during my internship warranted explicit
documentation because each involved a meaningful tradeoff between
competing engineering constraints. This section records the decision,
the alternatives considered, and the rationale for the chosen approach.

\subsection{Decision 1: Secrets Manager 64\,KB Limit \texorpdfstring{$\to$}{→} S3 Migration}
\label{sec:decision_secrets_s3}

\paragraph{Problem.}
The CloudFormation IaC provisioner Lambda function stores its
CloudFormation YAML templates as a payload retrieved at deploy time.
Initially, this payload was stored in AWS Secrets Manager as a single
JSON-encoded string. As the monitoring scope expanded to cover all seven
AWS service categories, the generated YAML grew to \textbf{14,720~lines}
at version v728 --- exceeding the \textbf{64\,KB hard limit} on
Secrets Manager secret values. The Lambda function began failing with
\texttt{SecretSizeExceeded} on the \texttt{PutSecretValue} call.

\paragraph{Options Considered.}

\begin{enumerate}
    \item \textbf{Split into multiple secrets}: Partition the YAML
          into $n$ secrets, each under 64\,KB, and have the Lambda
          function reassemble them. This preserves Secrets Manager
          as the storage layer but introduces assembly logic and
          a hard dependency on secret ordering.

    \item \textbf{Compress the payload}: gzip-compress the YAML before
          storing it as a binary secret. Achieves $\approx$70\%
          compression on YAML, bringing the payload well under 64\,KB.
          Adds a decompression step in the Lambda function.

    \item \textbf{Migrate to S3}: Store the YAML templates as S3 objects,
          versioned using S3 object versioning. The Lambda function
          retrieves templates via \texttt{s3:GetObject} at deploy time.
\end{enumerate}

\paragraph{Decision.} I chose \textbf{S3 migration}. The rationale:

\begin{itemize}
    \item No size limit --- templates can grow arbitrarily as additional
          AWS service categories are onboarded.
    \item Native versioning via S3 object versioning, providing a rollback
          mechanism for template changes at no additional cost.
    \item No reassembly or decompression complexity in the Lambda function.
    \item Cost negligible: a single \texttt{GetObject} call per deploy
          adds less than \$0.001 to operational cost.
\end{itemize}

\paragraph{Tradeoff.}
The Lambda execution role now requires \texttt{s3:GetObject} on the
template bucket, broadening its IAM surface slightly. The
\texttt{GetObject} call adds approximately 50~ms of latency to the
provisioner bootstrap --- acceptable for a deploy-time operation
that runs infrequently.

\subsection{Decision 2: AppName Tag as Mandatory Secondary Discriminator}
\label{sec:decision_appname}

\paragraph{Problem.}
As described in Section~\ref{sec:inventory_api}, the \texttt{CITCOWORKS}
and \texttt{MERIDIAN} environments share \texttt{BudgetCode=``Shared Cloud
Artifacts''}. A single-stage \texttt{BudgetCode} lookup cannot distinguish
resources from the two environments, causing inventory contamination.

\paragraph{Root Cause.}
Both platforms were onboarded before the IA team established the
policy requiring each platform to have a unique
\texttt{BudgetCode} tag value. Retroactively changing the tag across
all existing resources in two production environments would require
a coordinated tagging campaign with risk of operational disruption.

\paragraph{Decision.}
I designated the \texttt{AppName} tag as the \textbf{mandatory secondary
discriminator} for all resources in the \texttt{``Shared Cloud Artifacts''}
BudgetCode. Discoverers for \texttt{CITCOWORKS} and \texttt{MERIDIAN}
must implement the two-stage filter described in
Section~\ref{sec:inventory_api}.

\paragraph{Tradeoff.}
This decision creates a tagging governance obligation: all new resources
added to either environment \emph{must} carry the correct \texttt{AppName}
tag. A resource missing this tag will be invisible to the inventory API,
creating a silent coverage gap. This risk is mitigated by a CloudWatch
tag compliance alarm and documented in the \texttt{ia-sre-resources-inventory}
\texttt{README} as a platform-wide architectural constraint.

\subsection{Decision 3: CloudWatch as Sole Data Source (No Prometheus)}
\label{sec:decision_cloudwatch_only}

\paragraph{Problem.}
A production-grade Grafana deployment typically uses Prometheus as the
time-series database, with CloudWatch Exporter bridging AWS metrics.
This is the pattern recommended by the Grafana community for large-scale
AWS environments.

\paragraph{Alternative.}
Deploy a self-managed Prometheus server (on ECS or EC2) with the
CloudWatch Exporter sidecar, and configure Grafana to use Prometheus
as its data source.

\paragraph{Decision.} I used the \textbf{native Grafana CloudWatch plugin}
exclusively.

\paragraph{Rationale.}

\begin{itemize}
    \item \textbf{No additional infrastructure}: A Prometheus server
          and exporter fleet would introduce new ECS tasks, IAM roles,
          configuration management, and an additional failure domain.

    \item \textbf{Amazon Managed Grafana constraint}: AMG does not
          support self-managed Prometheus as a built-in data source
          without a Prometheus-compatible remote write endpoint.
          Connecting AMG to a self-managed Prometheus would require
          Amazon Managed Service for Prometheus (AMP), adding further
          cost and complexity.

    \item \textbf{Direct IAM authentication}: The native CloudWatch
          plugin authenticates directly via IAM roles, requiring zero
          credential management beyond the initial role configuration.

    \item \textbf{Cost}: CloudWatch metrics are already being published
          by AWS services at no marginal cost. Prometheus would require
          scraping, storage, and retention management.
\end{itemize}

\paragraph{Tradeoff.}
The accepted limitation is that CloudWatch Metric Insights SQL does not
support \texttt{p50}/\texttt{p99} percentiles in \texttt{GROUP BY} mode
(see Section~\ref{sec:grafana_amg_platform}). Latency distribution
dashboards must use \texttt{AVG} and \texttt{MAX} as proxies, or fall
back to the standard metrics wildcard mode which only shows recently
active functions. For the current monitoring maturity of the IA platform,
this tradeoff is acceptable --- the P50/P99 gap can be addressed in a
future phase by adding custom Lambda Duration metric emission with
explicit percentile dimensions.

\subsection{Decision 4: Lambda-Backed IaC Provisioner vs.\ Static CloudFormation}
\label{sec:decision_lambda_provisioner}

\paragraph{Problem.}
Provisioning one CloudWatch alarm per Lambda function for all monitored
metric types would generate over 100~\texttt{AWS::CloudWatch::Alarm}
resources per stack. Approaching the \textbf{CloudFormation 500-resource
per-stack limit} was a realistic concern as the fleet grew. A static
template encoding all alarms explicitly would also require manual
updates every time the Lambda fleet changed.

\paragraph{Alternative.}
Static CloudFormation templates with each alarm resource hardcoded ---
one YAML block per function per alarm type, updated manually when
functions are added or removed.

\paragraph{Decision.}
I chose a \textbf{Lambda-backed dynamic provisioner} that discovers the function
fleet at deploy time and synthesizes alarm resources programmatically.

\paragraph{Rationale.}

\begin{itemize}
    \item \textbf{Fleet-scale coverage}: The provisioner ensures every
          Lambda function in every platform has alarms
          configured, regardless of when it was deployed. No manual
          update is needed when a new function is added.

    \item \textbf{Avoids static YAML maintenance}: A static template
          for 119~functions $\times$ 3~alarm types $= 357$~resources
          would be a 500+ line YAML file requiring hand-editing for
          every fleet change.

    \item \textbf{Parameterization}: Thresholds, evaluation periods,
          and SNS topic ARNs are all externalised as SAM parameters,
          enabling multi-environment reuse without code changes.
\end{itemize}

\paragraph{Tradeoff.}
The Lambda provisioner introduces pipeline ordering complexity: the
provisioner Lambda must complete successfully before the CloudFormation
stack update can proceed, creating a sequential dependency. A Lambda
execution failure (timeout, permissions error) would halt the deployment
pipeline without deploying updated alarms. This is mitigated by
comprehensive error logging and a manual re-invocation pathway via
\texttt{python main.py deploy --phase 3}.

\begin{table}[ht]
\centering
\caption{Design Decision Summary}
\label{tab:design_decisions}
\begin{tabular}{p{3.2cm}p{3.5cm}p{5.0cm}}
\toprule
\textbf{Decision} & \textbf{Chosen Approach} & \textbf{Key Tradeoff} \\
\midrule
Template storage & S3 (from Secrets Manager) & +\texttt{s3:GetObject} permission; $\sim$50\,ms latency \\
\midrule
Tag disambiguation & AppName mandatory secondary key & Tags governance obligation; missing tag = invisible resource \\
\midrule
Metrics data source & CloudWatch native plugin only & No p50/p99 in GROUP BY mode; AVG/MAX as proxies \\
\midrule
Alarm provisioning & Dynamic Lambda provisioner & Sequential deploy dependency; Lambda failure halts pipeline \\
\bottomrule
\end{tabular}
\end{table}

Taken together, these four decisions shaped an architecture that favours
operational simplicity and zero-maintenance fleet coverage over strict
metric precision, while accepting bounded and well-understood tradeoffs
at each decision point. I documented the design in the respective
repository \texttt{README} and JIRA acceptance criteria so that future
team members can revisit any decision as the platform's monitoring
maturity evolves.



% ────────────────────────────────────────────────────────────
% PHASED DELIVERY
% ────────────────────────────────────────────────────────────
\section{Phase 1 --- Onboarding and CloudWatch (Weeks 1--7, Mar~15 -- May~10)}
\label{sec:phase1_onboarding}
% ────────────────────────────────────────────────────────────

\subsection{Part-Time Onboarding (Mar 15 -- Apr 12)}

My internship opened in part-time mode at roughly 10--20 hours per week, during
which my primary objective was orientation rather than delivery. I filed access requests
on Day~1 for AWS console, JIRA, Confluence, and CodeCommit; the approvals
trickled in over the first fortnight as each request was routed through the appropriate
security review. During this period, daily standups served as the most effective
source of context: listening to incident retros, deployment discussions, and sprint
reviews built a working mental model of the IA team's operations before I wrote any single
line of code.

Architecture documentation on Confluence gave a structural picture of the five
platforms (\texttt{automationhub}, \texttt{bpm}, \texttt{citcoworks},
\texttt{ia-bpo}, \texttt{meridian}, \texttt{rpacfs1}) and the Lambda, ECS, and EC2 AWS services they
contained. The recurring theme in these documents was a visibility gap: teams learned
about production failures from business users rather than from monitoring tools. That
pattern would become the guiding problem statement for my entire internship.

\subsection{Training Week (Apr 13--18)}

The week of April~13 was designated as a formal knowledge-transfer week.
I attended four focused sessions over three days:

\begin{itemize}
    \item \textbf{Ho Hoi Leung (Horace) --- RPA Support.}
          Horace walked through the production support rotation model: who is
          on-call each week, how the VPL queue opens at 5~PM~EST each weekday,
          how Blue Prism and UiPath Cloud Orchestrator are monitored, and
          the escalation ladder from Tier~1 business users through Tier~2 IA
          Operations to Tier~3 IA Development. This KT session established the
          mental framework that made the later support rotations tractable.

    \item \textbf{Bhanu Teja Murari --- CitcoWorks and Event Store.}
          This session covered the CitcoWorks platform architecture and, critically,
          the Event Store structured logging standard: every automation process emits
          events with a \texttt{subject} field matching the Automation Catalog project
          code, a \texttt{process\_identifier}, and a \texttt{run\_id} that threads
          together all log events for a single execution. Understanding this structure
          was essential for interpreting CloudWatch Logs Insights queries later in the
          Grafana build.

    \item \textbf{Meridian Team --- Kafka Event Pipeline.}
          The Meridian platform ingests events through a three-stage pipeline: the
          Event Routing Service receives events from source systems, the Action
          Processor applies business routing rules, and the OpenSearch Consumer
          persists events into the search index used by downstream services. This
          pipeline topology --- and its known sensitivity to consumer lag --- would
          directly inform the MSK and OpenSearch monitoring rows of the Grafana
          dashboard.

    \item \textbf{\AE xeo --- Treasury Platform.}
          A brief session on the \AE xeo cloud-based wire approval and fund movement
          platform, which several VPL automation workflows interact with via HTTP
          and UI automation.
\end{itemize}

The training week also included an unplanned but instructive event. On
April~17, the Meridian platform suffered a significant processing delay: a surge
of automated events from Citco Recs congested the OpenSearch consumption layer,
causing approximately 9~million Kafka messages to queue while downstream services
continued polling for completions that never arrived. The Winton and GS fund teams
flagged the issue to Dublin support around 8~AM; root cause was not identified until
9:30~AM and full restoration took until approximately noon---a 7.5-hour end-to-end
delay. Observing the incident response process firsthand made abstract discussions
about ``observability gaps'' concrete: the platform had no alerting on
\texttt{EstimatedMaxTimeLag} for any Kafka consumer group, and the team had no
single-pane view of Meridian health. Both gaps were named explicitly as motivators
for IISS-487 in the incident retrospective.

\subsection{First Deliverables --- IISS-455, IISS-461, IISS-469 (Apr 19 -- May 10)}

\subsubsection{IISS-455: Lambda Dashboard and Inventory (Apr 7)}

The first concrete deliverable was a CloudWatch dashboard covering the
\texttt{rpacfs1} and \texttt{ia-bpo} Lambda function fleets, which I completed in
approximately two hours on April~7. I provisioned the dashboard via
CloudFormation and used \texttt{SEARCH()} expressions to enumerate Lambda
functions by \texttt{FunctionName} prefix, providing per-function panels for
invocation counts, error rates, and duration averages.

Simultaneously, I produced a Confluence inventory page (IISS-461) mapping the
\textbf{51 Lambda functions} then visible across the two platforms to their
owning team and purpose. This manual enumeration was my first attempt to answer
a question the team had been unable to answer from the console alone: ``how many
Lambda functions do we actually have?''

\subsubsection{IISS-469: Extended Dashboard and Architecture Analysis (Apr 15--21)}

IISS-469 extended the dashboard to cover the \texttt{cfsbpm}, \texttt{citcoworks},
and \texttt{meridian} platforms. The extension work surfaced two structural
limitations in the \texttt{SEARCH()}-based approach that would permanently shape
my Grafana architecture:

\begin{itemize}
    \item \textbf{SEARCH() visibility gap.} The \texttt{SEARCH()} function only
          returns metrics for resources with recent data points. A Lambda function
          that has not been invoked within the lookback window disappears from the
          dashboard entirely, creating a misleading picture of the fleet size.
          This gap I confirmed empirically: re-querying after a quiet weekend
          returned 38 functions where Monday morning returned 51.

    \item \textbf{Top-K and metric filter constraint.} A proposed central
          log-aggregation design --- routing all Lambda function log events into
          a shared subscriber --- would require \textbf{106 metric subscription
          filters} to cover all function--metric-dimension combinations across the
          six platforms. This was blocked by an IAM permission boundary:
          \texttt{logs:PutSubscriptionFilter} was not granted within the IA team's
          scoped permission set.
\end{itemize}

I also produced a documented analysis of five approaches to out-of-memory
(OOM) detection for Lambda functions: (1)~parsing \texttt{REPORT} log lines via
Logs Insights; (2)~comparing \texttt{MaxMemoryUsed} against \texttt{MemorySize}
via the Lambda enhanced monitoring extension; (3)~CloudWatch Lambda Insights
(paid tier); (4)~custom metric emission on threshold breach; and
(5)~structured log parsing via metric filter patterns. I scored each option
on detection latency, cost, and permission requirements. This analysis fed
directly into my alarm design under IISS-487.

By the end of Phase~1, the full Lambda inventory across all six platforms was documented
across all six platforms, documented in Confluence. The \texttt{SEARCH()}
visibility gap and the IAM filter constraint had established a clear design
directive: my next generation monitoring platform would need to use CloudWatch
Metric Insights SQL rather than metric subscription filters.

% ────────────────────────────────────────────────────────────
\section{Phase 2 --- Grafana Intensive (Weeks 8--11, May~6 -- May~27)}
\label{sec:phase2_grafana}
% ────────────────────────────────────────────────────────────

\subsection{Elevation to Critical and Week~8 Support Shift}

On \textbf{May~6}, Kishor elevated IISS-487 to Critical priority, formalising
what had been understood informally since the April~17 incident: the IA team
needed a centralized, real-time monitoring platform before the next Meridian
outage occurred. The elevation came during my first support rotation week
(Week~8), which meant the Grafana work could not begin in earnest until the
following Monday. The support shift provided an early orientation to the
operational rhythms that the dashboard was being built to serve.

\subsection{Infrastructure Deployment Day --- May 11}

May~11 was the single most infrastructure-dense day of my internship. I deployed four
CloudFormation stacks in sequence:

\begin{enumerate}
    \item \textbf{\texttt{ia-sre-lambda-alarms}} --- CloudWatch alarms for Lambda
          error rates, throttle counts, and duration percentiles across all five
          budget codes.
    \item \textbf{\texttt{ia-sre-ecs-alarms}} --- ECS CPU and task-failure alarms,
          enabled after Standard Container Insights was activated on five of the
          six Meridian ECS clusters.
    \item \textbf{\texttt{ia-sre-sqs-alarms}} --- Queue depth and
          \texttt{ApproximateAgeOfOldestMessage} alarms per monitored queue.
    \item \textbf{\texttt{ia-sre-msk-alarms}} --- MSK consumer lag alarm on
          \texttt{EstimatedMaxTimeLag}, the metric whose absence had made the
          April~17 incident invisible until business users flagged it.
\end{enumerate}

Alongside the alarms, I provisioned the Amazon Managed Grafana (AMG) workspace
and configured the 16-role IAM agent permission model to give the Grafana
service account cross-account read access across all IA environments. IAM
permission boundaries were the main obstacle of the day: the
\texttt{cloudwatch:GetMetricData} and \texttt{logs:StartQuery} permissions
required a non-trivial negotiation with the scoped permission set governing
the internship AWS account. I enabled Container Insights on five of the six
Meridian clusters, beginning to populate the \texttt{ECS/ContainerInsights}
namespace with task-level CPU and memory metrics.

\subsection{Grafana 10.4 to 12.4 Upgrade --- May 13}

Two days after the initial infrastructure deployment, the AMG workspace was
upgraded from Grafana \textbf{10.4 to 12.4}. This was unplanned and introduced
breaking changes in two areas:

\begin{itemize}
    \item \textbf{CloudWatch Metric Insights SQL editor.} The query syntax
          editor in Grafana~12.4 changed how SQL expressions were associated
          with panel targets. Queries that I had validated against the 10.4
          editor were no longer recognised in the new model.
    \item \textbf{Variable interpolation.} Template variable syntax changed
          between versions, causing several \texttt{\$budget\_code} and
          \texttt{\$ecs\_cluster} references to resolve incorrectly --- panels
          that had returned data were now blank.
\end{itemize}

My recovery required a systematic audit of all panels built in the first two days,
identifying which queries used version-specific syntax and migrating them to the
stable Grafana~12.4 CloudWatch query model. Approximately one working day was
absorbed by this migration. The experience directly motivated NFR6 --- documented
in the design as ``use only standard CloudWatch plugin queries; avoid
AMG-version-specific syntax'' --- and enforced a panel-testing discipline:
I verified every new panel in the browser before committing it.

\subsection{Dashboard Build --- May 11--27 (80+ Hours)}

The three-week block from May~11 to May~27 was the most concentrated development
period of my internship, logging over 80~hours on IISS-487 alone. I organised progress
around implementing the dashboard rows in sequence, with cross-cutting
infrastructure problems resolved as they arose.

\textbf{Rows 1--3 (SQS and Lambda), May~11--15.} The SQS row required a
significant query architecture decision. My initial approach used a Grafana
\texttt{\$sqs\_queue} template variable to drive a per-queue metric panel, but
this produced empty data for idle queues --- the same SEARCH() visibility gap
encountered in Phase~1 now appearing in Grafana. I solved this by switching to
Metric Insights SQL with \texttt{GROUP BY QueueName}, returning all queues as
separate time series within a single query. Lambda error-rate panels were
initially slow to load because each panel made two CloudWatch API calls
(Errors and Invocations) to compute the ratio in Grafana; I replaced these
with pre-computed error rate expressions evaluated in Metric Insights SQL,
significantly reducing panel load time.

\textbf{ECS Row --- SQL Migration and the \texttt{\$platform} Variable, May~14--15.}
The initial ECS panels used basic CloudWatch metric queries, which did not
support per-service breakdown across 22 clusters. I migrated these to
Metric Insights SQL with \texttt{GROUP BY ServiceName}, enabling a single
panel to render per-service CPU or memory utilisation as distinct time series.
I introduced a dashboard template variable \texttt{\$platform} so that ECS
panels could be filtered dynamically by environment, matching the pattern
already in use for the SQS and Lambda rows. The initial SQL used
\texttt{SCHEMA()} syntax, which hit CloudWatch query limits; I replaced this
with the \texttt{FROM "ECS/ContainerInsights"} namespace form.

\textbf{Percentile Limitation, May~14.} A key CloudWatch constraint was
confirmed during my Lambda Duration panel work: Metric Insights SQL \texttt{GROUP BY}
mode supports only \texttt{AVG}, \texttt{SUM}, \texttt{MIN}, and \texttt{MAX}.
Percentile statistics (\texttt{p50}, \texttt{p99}) are unavailable in this mode.
For the Lambda duration P50/P95 panels, I implemented a wildcard dimension workaround
--- querying the standard CloudWatch metrics API with
\texttt{FunctionName:~["*"]} --- accepting the visibility gap tradeoff for
recently inactive functions.

\textbf{MSK, DocumentDB, OpenSearch Rows, May~19--22.} I built the MSK/Kafka row
around the \texttt{EstimatedMaxTimeLag} metric per consumer group,
directly implementing the alarm that would have detected the April~17 incident.
DocumentDB and OpenSearch rows followed, with the OpenSearch JVM heap panel
receiving particular attention given Meridian's dependency on OpenSearch for
event querying.

\textbf{Secrets Manager 64~KB Limit and S3 Migration, May~21.}
As the dashboard grew, the Lambda-backed CloudFormation provisioner's payload
--- stored as a JSON-encoded string in AWS Secrets Manager --- exceeded the
service's \textbf{64~KB hard limit}. By version~v438 of the provisioner, the
assembled CloudFormation YAML had grown to \textbf{14,720~lines}. The provisioner
Lambda began failing with \texttt{SecretSizeExceeded} on the
\texttt{PutSecretValue} call. I migrated the templates to S3 object storage
on the same day, updating the provisioner Lambda to read templates via
\texttt{s3:GetObject} at deploy time. This decision is discussed in depth in
Section~\ref{sec:decision_secrets_s3} of the architecture chapter.

\textbf{Alarm Notifications and Composite Alarms, May~27--29.}
I deployed a Re-Notifier Lambda to ensure alarm state changes were routed
through an SNS topic to the IA-SRE email distribution list, regardless of
whether CloudWatch had already sent a prior notification in the same alarm
state. I configured eight composite alarms to aggregate per-service alarm
states into a single \texttt{OK}/\texttt{ALARM} indicator per budget code
environment, displayed on Row~12 of the dashboard.

\textbf{Phase~2 Milestone.} By May~27, my core monitoring platform was
functional: the MSK consumer lag alarm was active, the first six dashboard rows
were complete, and the composite alarm structure was in place. The detection
latency for a Meridian Kafka congestion event had been reduced from the observed
3.5-hour business-team lag of April~17 to a projected sub-30-second alarm fire
time.

% ────────────────────────────────────────────────────────────
\section{Phase 3 --- Dashboard Completion and SQS Feature (Weeks 12--15, May~28 -- Jun~18)}
\label{sec:phase3_completion}
% ────────────────────────────────────────────────────────────

\subsection{Rows 7--14 and the Deep-Link Pattern (May~28 -- Jun~4)}

With the first six rows stable, Phase~3 completed the remaining dashboard structure.
Rows~7--8 added EC2 monitoring using the \texttt{AppName} tag as the primary
discovery dimension --- the fallback pattern I established in IISS-483 to handle
the \texttt{rpacfs1} tag mismatch (where \texttt{BudgetCode} carried the value
\texttt{"CFS/RPA"} rather than \texttt{"rpacfs1"}). The EC2 row covered
\textbf{83 instances} across all environments, with per-instance CPU, memory
(via CloudWatch Agent), and network panels.

Rows~9--10 added DocumentDB and OpenSearch monitoring. I gave the OpenSearch panels ---
index rate, search latency, and JVM heap pressure --- priority attention
because OpenSearch was the bottleneck in the April~17 incident. I wired a JVM heap
utilisation panel above 75\% to a CloudWatch alarm to provide early
warning before the service degraded under load.

\textbf{Deep-link panels (Row~14)} were a usability feature not originally in
scope but I added them after a discussion with Kishor: each resource panel now includes
a clickable URL that navigates directly to the relevant AWS Console page
(Lambda function, ECS service, EC2 instance, or MSK cluster) without requiring
the operator to re-navigate the console hierarchy during an incident. These
panels used Grafana's Text panel with an HTML data link, deriving the console
URL from the resource ARN or dimension value already present in the panel data.

\subsection{IISS-482: SQS Feature --- Jun 12}

The SQS monitoring feature (IISS-482) I implemented on June~12 was a
self-contained sub-deliverable. I added a \textbf{registry-based
SQS monitoring pipeline} with a three-tier fallback model for resilient queue
data retrieval:

\begin{enumerate}
    \item \textbf{Tier~1 --- Live API.} Query the CloudWatch \texttt{AWS/SQS}
          namespace directly for real-time queue depth and age metrics.
    \item \textbf{Tier~2 --- S3 Cache.} If the live API fails, retrieve the
          last-known queue snapshot from an S3 bucket. Stale by up to one cache
          TTL window, but always available.
    \item \textbf{Tier~3 --- Local File.} If both API and S3 fail, fall back to
          a locally bundled queue inventory file reflecting the deployment-time
          state of the fleet.
\end{enumerate}

I accompanied the SQS feature with \textbf{960 automated tests} covering registry
dispatch logic, fallback transitions, queue metadata serialisation, and the
dashboard builder integration. Writing the test suite before the feature code
was complete surfaced two edge cases in the fallback transition logic --- a
race condition in the Tier~1 to Tier~2 failover path, and an incorrect
queue-name normalisation step --- both of which I fixed before the feature
was merged.

\subsection{Final Deployment --- Jun 18, Version v728}

The dashboard reached its production-ready state on June~18 with the deployment
of version \textbf{v728}. The final deployment session logged 19 hours across
multiple sessions that day, covering the last round of deep-link fixes,
DocumentDB and Redis restructuring, and the release merge to the main branch.
My cumulative totals at release were:

\begin{itemize}
    \item \textbf{140 hours} logged on IISS-487
    \item \textbf{14 dashboard rows}, \textbf{110+ panels}
    \item \textbf{38 CloudWatch alarms} active across Lambda, ECS, MSK, SQS, and EC2
    \item \textbf{7 AWS services} monitored under a single Grafana workspace
    \item \textbf{4 CloudFormation stacks} managing the alarm infrastructure
\end{itemize}

% ────────────────────────────────────────────────────────────
\section{Phase 4 --- Inventory API and Refinements (Weeks 16--22, Jun~19 -- Aug~31)}
\label{sec:phase4_inventory}
% ────────────────────────────────────────────────────────────

\subsection{IISS-507: Scaffolding the Inventory API --- Jun 11}

Work on IISS-507 began on June~11 --- one week before IISS-487 was formally
closed --- taking advantage of a brief quieter period during the dashboard
polish phase. I housed the project in the \texttt{ia-sre-resources-inventory}
repository and scaffolded it using Kiro IDE's spec-driven development workflow:
I wrote a \texttt{requirements.md} first, from which Kiro's sub-agents
generated a \texttt{design.md} and \texttt{tasks.md}, and then produced stub
implementations for each task in the plan. The spec-first approach reduced my
scaffolding time to approximately 2.5~hours for a project that would otherwise
have required a full day of boilerplate setup.

I chose the technology stack to align with the team's existing Python
ecosystem: \texttt{flask[async]} for route handlers,
\texttt{boto3} for all AWS resource discovery, Pydantic for response model
validation, and \texttt{serverless-wsgi} for Lambda deployment.

\subsection{Inventory API Development --- Jun 11 -- Jul 14}

\subsubsection{Lambda Discoverers}

I implemented Lambda discoverers for all six platforms, confirming the
Lambda inventory documented in Phase~1. The \texttt{rpacfs1} discoverer
used the \texttt{AppName} fallback pattern: a primary filter on
\texttt{BudgetCode="rpacfs1"} returned zero results, triggering the fallback
to \texttt{AppName="rpacfs1"}, which correctly identified the fleet. This
fallback became the standard for all \texttt{rpacfs1} resource types.

\subsubsection{ECS Discoverers and Tag Disambiguation}

I implemented ECS discoverers for all 22 clusters across six platforms.
The \texttt{CITCOWORKS} and \texttt{MERIDIAN} environments presented the most
complex discovery challenge: both share the AWS tag value
\texttt{"Shared Cloud Artifacts"} for the \texttt{BudgetCode} dimension, a
legacy of their original onboarding before the IA team established per-platform
uniqueness requirements. The solution was a two-stage filter in each discoverer:
Stage~1 filters by \texttt{BudgetCode="Shared Cloud Artifacts"}; Stage~2
applies \texttt{AppName="meridian"} or \texttt{AppName="citcoworks"} to
discriminate between the two environments. I formally designated the \texttt{AppName} tag
as the mandatory secondary discriminator for all resources
in the shared-tag platforms, a governance decision documented in
the repository \texttt{README} and enforced by a tag-compliance CloudWatch alarm.

\subsubsection{EC2 Discoverers}

My EC2 discoverers covered \textbf{83 instances} with the same AppName fallback
pattern used in the Lambda discoverers. The Pydantic response model for each
discovered resource carried five fields: \texttt{name}, \texttt{type},
\texttt{budget\_code}, \texttt{tags} (as a free-form dictionary), and
\texttt{arn} for direct AWS Console linkage.

\subsubsection{Async Architecture and API Latency}

I designed the route handlers to use \texttt{asyncio.gather()} to dispatch all discoverers for
a given resource type concurrently rather than sequentially. For a
\texttt{GET /api/v1/lambda} request without a budget-code filter, this reduces
the response time from the sequential sum of five discoverer latencies to
approximately the latency of the single slowest AWS API call --- consistently
keeping the endpoint within the 10-second NFR target across the full Lambda fleet.

\subsection{Test Suite Milestone --- Jul 14: 562 Tests at 99\% Coverage}

By July~14, my test suite reached \textbf{562 automated tests} with
\textbf{99\% code coverage}. The suite comprised four layers:

\begin{itemize}
    \item \textbf{Unit tests} --- each discoverer tested in isolation with mocked
          boto3 clients, verifying correct tag filtering, pagination handling, and
          \texttt{Resource} object construction.
    \item \textbf{Integration tests} --- full Flask test client hitting all REST
          endpoints, verifying routing, serialisation, and error response format.
    \item \textbf{Property-based tests (Hypothesis)} --- tag-parsing logic and
          registry dispatch tested with randomly generated inputs, verifying that
          the two-stage filter never returns a \texttt{MERIDIAN} resource when
          \texttt{budget\_code=citcoworks} is specified, regardless of tag value
          combinations.
    \item \textbf{Edge case tests} --- malformed tags, empty platforms,
          boto3 pagination edge cases (empty pages, single-page results, maximum
          page sizes).
\end{itemize}

\subsection{X-Ray Tracing Research --- Jul--Aug}

Following the dashboard completion, my research attention turned to IISS-489's
roadmap item for distributed tracing. The X-Ray analysis evaluated two
instrumentation approaches: the legacy \texttt{aws-xray-sdk} and the newer
\texttt{aws-lambda-powertools[tracer]}, with the latter preferred for new Python
Lambda functions due to its decorator-based API and tighter integration with the
PowerTools structured logging layer already in use.

My cost modelling estimated approximately 5--10~ms of per-invocation overhead
from X-Ray SDK calls, with sampling rates configurable between 1\% and 100\%
to control cost. At the current Lambda invocation volume, a 5\% sampling rate
was projected to add less than \$2/month to the CloudWatch spend while providing
statistically representative trace coverage for latency anomaly detection.

I ultimately deferred X-Ray deployment to the post-internship roadmap (IISS-837
through IISS-844) for two reasons: the IAM role configuration for
\texttt{xray:PutTraceSegments} required a cross-team permission review, and the
remaining internship weeks were more effectively spent on documentation and
knowledge transfer to the permanent team.

I dedicated the final weeks of August to runbook updates, Confluence documentation
for the dashboard operating model, and structured knowledge transfer sessions with
Horace and Kri. The Mitacs BSI renewal for the September~1 -- December~31, 2026
extension was confirmed, providing continuity for the X-Ray and log-group discovery
work scoped in the Phase~4 JIRA backlog.

% ────────────────────────────────────────────────────────────
\section{RPA Support Operations}
\label{sec:rpa_support}
% ────────────────────────────────────────────────────────────

\subsection{Support Rotation Model}

Four dedicated support rotation weeks (Weeks~8, 12, 16, and~20) ran in parallel
with my primary development phases. During a support week, my primary
responsibility shifted from development to RPA production support: monitoring
Blue Prism and UiPath bots, responding to service desk tickets,
performing configuration changes, and escalating or resolving incidents within
the shift window. The operational rhythm of each support day was structured
around the following checkpoints:

\begin{itemize}
    \item \textbf{VPL queue monitoring} --- the VPL queue opens at 5~PM~EST
          each weekday; checking bot health, queue depth, and processing status
          at queue open was the first task of each evening shift.
    \item \textbf{Corporate Actions triage} --- three process tracks
          (Mandatory, Voluntary, and Ops Transfer) each have configuration files
          on BPShare that must be maintained for new client onboarding.
    \item \textbf{Blue Prism Orchestrator health checks} --- robot session status,
          failed jobs, and any processes stuck in ``running'' state.
    \item \textbf{UiPath Cloud Orchestrator health} --- robot heartbeat status,
          stuck or failed jobs in the COE/HedgeFunds/PortfolioDataops folders.
    \item \textbf{On-call escalations} --- business team contacts from
          Dublin, London, and Halifax flagging processing failures via
          service desk tickets or direct chat.
\end{itemize}

I handled roughly a dozen service desk tickets across the four rotation
weeks, covering client onboarding (Corporate Actions FileMasks updates),
process retriggering (MESO and FX Closeout), log investigation, and UI
selector troubleshooting. Four incidents rose to the level of JIRA tracking
and are described below.

\subsection{Notable Incidents}

\subsubsection{IISS-706 --- CAIS Deadline Hotfix (Jun~8, 7 Hours)}

At approximately 9~AM on June~8, the CAIS Deadline process failed with a
cryptic \texttt{SMBAuthenticationError: SpnegoError} when uploading output
to the UCM file share. The error message --- \textit{``No username or password
was specified and the credential cache did not exist or contained no
credentials''} --- provided no direct indication of the root cause.

My investigation path proceeded as follows: the CloudWatch logs showed a
\texttt{SpnegoError} deep inside the \texttt{smbprotocol} library;
the SMB protocol uses NTLM authentication, which requires a valid
username--password pair; the credentials are loaded from AWS Secrets Manager
by \texttt{constants.py} at module import time; checking the actual secret
values revealed that the most recent Secrets Manager rotation had left the
\texttt{ucm\_username} and \texttt{ucm\_password} fields as empty strings
during the rotation window. Because \texttt{constants.py} caches the credential
values as module-level variables for the entire container lifecycle, a container
that started during that empty-credential window continued to use empty
credentials for every subsequent UCM upload --- until an ECS task restart
triggered a fresh container with a fresh module import.

This explained the intermittent pattern observed in the preceding days: the
process had failed on June~3 and succeeded on June~4 (when the container had
been restarted) and failed again on June~8 (when a new container picked up
empty credentials during another rotation window).

My hotfix implemented two layers of defence:

\begin{enumerate}
    \item \textbf{Fail-fast credential validation at startup.} In
          \texttt{constants.py}, after loading the secret, an explicit check
          raises a \texttt{RuntimeError} immediately if either credential field
          is empty or absent. The container fails loudly at startup rather
          than silently until the first UCM call.

    \item \textbf{Call-time credential pre-check.} In \texttt{pysmb\_helper.py},
          a \texttt{ValueError} is raised at the start of
          \texttt{upload\_to\_share()} and \texttt{download\_from\_share()} if
          the \texttt{user} or \texttt{passw} arguments are empty. This provides
          a second line of defence if the credential validation in
          \texttt{constants.py} is bypassed by an alternate import path.
\end{enumerate}

My hotfix was reviewed, approved, and merged to the release branch on the same
day. The CAIS filing completed before the regulatory deadline. I documented the root cause ---
Secrets Manager rotation leaving a zero-credential window --- as
a permanent fix candidate requiring either periodic credential re-fetching
(lazy loading on each UCM call) or a rotation configuration change to ensure
the old secret remains valid until the new one is confirmed. This incident
directly motivated the UCM credential validation pattern I documented in
\texttt{constants.py} and the ECS SaturationAlarm I subsequently added to the
IISS-487 alarm set.

\subsubsection{IISS-538 --- MESO Autoscale Issue (Jun~5)}

The \texttt{ia\_meso\_process} (Middle Office Month-End Sign-Off) ECS service
was consuming excessive CPU and memory, with tasks scaling to the maximum
task count without shedding load proportionally. The root cause was an
ECS service autoscaling policy that had not been configured: without a scaling
policy, the service could spin up additional tasks but had no rule to scale
back down or to reject new tasks once the service was saturated. The fix was
a corrected ECS service autoscaling configuration. The incident motivated my
addition of an ECS CPU and memory \texttt{SaturationAlarm} to the IISS-487
alarm set, targeting any service with CPU sustained above 80\% for more than
five minutes.

\subsubsection{IISS-741 --- MESO Month-End Stuck Tasks (Jul~1, Critical)}

On July~1, INNOCAP and SAMLMOS fund tasks were found stuck in the queue with
no heartbeat from the executing Lambda. My manual log investigation confirmed that
the tasks had stalled and were not progressing. I resolved this by manually
retriggering the \texttt{ia\_meso\_process} pipeline for both affected funds,
following a KT session with Horace on the admin access required to mark stuck
tasks as failed. Processing completed approximately two hours after the initial
escalation, with no data loss. The incident reinforced the importance of
Lambda heartbeat monitoring --- a gap I subsequently noted in the IISS-487
roadmap as a candidate for the Phase~5 X-Ray tracing work.

\subsubsection{IISS-753 --- VPL IPV UI Selector Failure (Jul~8)}

The VPL Integrated Pricing and Valuation (\AE xeo) process encountered
repeated \texttt{CheckpointException} failures in the UiPath
\texttt{OpenValuationAnalysis.xaml} workflow --- specifically in the
``Retry Scope~--- Wait for Pricing Window'' step. The exception indicated that
the bot was waiting for the \AE xeo Pricing/Valuation window to become available
and exhausting its retry budget before the window was ready. Dynamic UI elements
in the fund system caused selector brittleness when the window load time varied
across sessions. I estimated the ticket at \textbf{16 hours} for a full
selector resilience investigation and escalated it to the development team
lead (Clarisa/Pennylaine) due to the Bitbucket repo access restriction and
the absence of updated exception-handling documentation. I identified this pattern --- UiPath
selector instability on dynamically loading AExeo screens --- as a recurring operational risk across multiple VPL workflows.

\subsubsection{IISS-876 --- VPL Pricing GTL Incorrect Booking (Jul~30)}

The VPL Pricing automation was booking incorrect values for GTL
(Guaranteed Transfer of Losses) funds. My root cause analysis identified a
targeted process logic error affecting how the GTL booking entry was calculated.
Resolution required my coordination with the business team to identify affected
NAV calculations and verify the corrected output. The incident highlighted the
importance of output reconciliation validation in financial automation ---
a process that checks calculated values against an independent source before
the booking step is executed.

\subsection{Support Incident Summary}

Table~\ref{tab:support_incidents} summarises the four notable support incidents
and their operational characteristics.

\begin{table}[ht]
\centering
\caption{Support Incident Summary}
\label{tab:support_incidents}
\begin{tabular}{p{1.6cm}p{1.5cm}p{4.0cm}p{4.5cm}p{1.0cm}}
\toprule
\textbf{JIRA} & \textbf{Date} & \textbf{Description} & \textbf{Resolution} & \textbf{Hours} \\
\midrule
IISS-706 & Jun 8 & CAIS Deadline \texttt{SpnegoError} --- empty UCM credentials from Secrets Manager rotation & Fail-fast credential validation in \texttt{constants.py} and \texttt{pysmb\_helper.py}; ECS restart & 7h \\
\midrule
IISS-538 & Jun 5 & MESO \texttt{ia\_meso\_process} ECS tasks consuming excessive CPU/memory, not autoscaling & ECS service autoscaling policy corrected; ECS SaturationAlarm added to IISS-487 & $\sim$2h \\
\midrule
IISS-741 & Jul 1 & INNOCAP/SAMLMOS MESO fund tasks stuck in queue, no Lambda heartbeat & Manual retrigger of pipeline for affected funds after admin access granted & $\sim$3h \\
\midrule
IISS-753 & Jul 8 & VPL IPV \AE xeo pricing window selector failure in UiPath (\texttt{CheckpointException}) & Escalated to dev team; 16h selector resilience investigation scoped & 16h est. \\
\midrule
IISS-876 & Jul 30 & VPL Pricing GTL incorrect booking values & Targeted process logic fix; business team reconciled affected NAV calculations & $\sim$4h \\
\bottomrule
\end{tabular}
\end{table}

\subsection{Cross-Cutting Learning: How Support Incidents Shaped Dashboard Design}

The four support incidents were not isolated events; each one contributed a
concrete requirement or architectural constraint to my Grafana dashboard and
the inventory API:

\begin{itemize}
    \item \textbf{IISS-538 (MESO autoscale)} motivated my addition of an ECS
          CPU and memory \texttt{SaturationAlarm} to the IISS-487 alarm set.
          Without the alarm, a repeat of the unbounded CPU growth would again
          require a business team escalation before the SRE team was aware.

    \item \textbf{IISS-706 (CAIS credentials)} surfaced the \texttt{constants.py}
          module-import-time credential pattern as a systemic fragility. The
          lesson --- that credential loading must account for Secrets Manager
          rotation windows --- was generalised into a design principle I recorded
          in the project runbook: any ECS process that uses long-lived
          module-level credential variables should include a container-startup
          health check that fails fast if credentials are empty.

    \item \textbf{IISS-741 (MESO stuck tasks)} demonstrated the value of a
          Lambda heartbeat or execution timeout alarm. If the executing Lambda
          had emitted a heartbeat metric every 30~seconds and an alarm had
          fired on 5~minutes of missing heartbeat, the stuck tasks would have
          been detected automatically rather than by manual queue inspection.
          I added this to the IISS-487 roadmap as a Phase~5 alarm type.

    \item \textbf{IISS-753 (VPL selector failure)} underlined UiPath selector
          brittleness as a recurring risk across all UI-based automations.
          While this falls outside the SRE monitoring scope, it informed my
          recommendation in the runbook that UI-based processes targeting
          \AE xeo should implement anchor-based selectors with a fallback path
          and a configurable retry window rather than fixed timeouts.
\end{itemize}

The support rotation model itself --- four dedicated weeks spread across my
internship --- was effective at preventing the two responsibilities from fully
competing. During non-support weeks, development could proceed without context
switches to service desk tickets. During support weeks, full operational
immersion built the domain knowledge (queue names, process flows, escalation
paths) that made my monitoring design decisions credible rather than abstract.
The cost was four weeks of development velocity; the benefit was a dashboard
that reflects the operational realities of the IA platform rather than an
idealised architecture.
